\section{Proof of the Affine Delayed Markov Bound}
\label{app:affine}
Fix a stationary action $a=(q,b,\eta)$ and condition on the history immediately before the retained joint observation.
All expectations below include the Markov observations; the delay profile is deterministic as specified in Assumption~\ref{ass:catalogue}.
Write $\bar H_k=q^{-1}\sum_iH_{i,k}$, $\bar\xi_k=q^{-1}\sum_i\xi_{i,k}$, and
\begin{equation}
v_k=e_k-\eta(\bar H_ke_k-\bar\xi_k),
\quad
r_k=\frac1q\sum_iH_{i,k}(e_{k-\tau_i}-e_k).
\label{eq:appendix-vr}
\end{equation}
Then Eq.~\eqref{eq:affine-td} is $e_{k+1}=v_k-\eta r_k$.

Total-variation duality and stationary centering give
\begin{equation}
\left\|\E[\bar\xi_k\mid\F_k]\right\|\le2G\delta.
\label{eq:conditional-bias}
\end{equation}
Moreover, stationarity and the bounded integrand imply
\begin{align}
\E[\|\bar H_ke-\bar\xi_k\|^2\mid\F_k]
&\le(2K_q+4L^2\delta)\|e\|^2\nonumber\\
&\quad+2\Omega_q+4G^2\delta.
\label{eq:second-moment-transfer}
\end{align}
Combining Eq.~\eqref{eq:conditional-bias} with the strong monotonicity of $A$ and applying Young's inequality with parameter $\mu_\delta$ yields
\begin{equation}
\E\|v_k\|^2\le a_\delta(\eta)R_k+\beta_\delta(\eta).
\label{eq:v-bound}
\end{equation}

The affine telescoping identity retains the innovation:
\begin{equation}
e_{k-\tau_i}-e_k=\eta\sum_{s=k-\tau_i}^{k-1}\left[\frac1q\sum_jH_{j,s}e_{s-\tau_j}-\bar\xi_s\right].
\end{equation}
Jensen's inequality and the almost-sure bounds in Assumption~\ref{ass:catalogue} give
\begin{equation}
\E\|r_k\|^2\le2\eta^2L^4\tau_{\rm rms}^2R_k+2\eta^2L^2G^2\tau_{\rm rms}^2.
\label{eq:r-bound}
\end{equation}
For any $\lambda>0$,
\begin{equation}
\|v_k-\eta r_k\|^2\le(1+\lambda)\|v_k\|^2+(1+\lambda^{-1})\eta^2\|r_k\|^2.
\end{equation}
Substituting Eq.~\eqref{eq:v-bound} and Eq.~\eqref{eq:r-bound} gives $\E\|e_{k+1}\|^2\le c_\lambda R_k+d_\lambda$.
The choice $\lambda=\sqrt{h_\tau/a_\delta}$ minimizes the contraction coefficient and produces Eq.~\eqref{eq:affine-cd}.
For zero delay, $r_k=0$ and the continuous limiting definitions follow directly.

Let $R_\star=d_{\rm aff}/(1-c_{\rm aff})$.
If $R_k\ge R_\star$, every new value entering the sliding envelope is at most $c_{\rm aff}R_k+d_{\rm aff}\le R_k$.
If $R_k<R_\star$, the same one-step bound keeps every new value below $R_\star$.
Thus, in both cases, each new value is at most $R_\star+c_{\rm aff}(R_k-R_\star)_+$.
After $2D+1$ updates the entire old window is replaced, giving $(R_{k+2D+1}-R_\star)_+\le c_{\rm aff}(R_k-R_\star)_+$.
Iteration gives $R_{nm}\le R_\star+c_{\rm aff}^n(R_0-R_\star)_+\le c_{\rm aff}^nR_0+R_\star$, proving Theorem~\ref{thm:affine}.

\subsection{Proof of Corollary~\ref{cor:robust-certificate}}
Write $\delta^+=C_q(r_q^+)^b$.
On the coverage event, $\delta^+\ge\delta$ and monotonicity on $\mu-2L\delta^+>0$ gives $a^+\ge a$ and $\beta^+\ge\beta$.
Evaluating the delayed Young inequality at the shared $\lambda^+=\sqrt{h/a^+}$ proves the upper recursion; separately optimized forcing terms need not be monotone.
For $\Delta\delta=\delta^+-\delta$, $\Delta K=K_q^+-K_q$, and $\Delta\Omega=\Omega_q^+-\Omega_q$,
\begin{align}
\Delta a
&=2\eta^2\Delta K+(2\eta L+4\eta^2L^2)\Delta\delta,\nonumber\\
\Delta\beta
&=2\eta^2\Delta\Omega+4\eta^2G^2\Delta\delta\nonumber\\
&\quad+4\eta G^2\left[
\frac{(\delta^+)^2}{\mu-2L\delta^+}
-\frac{\delta^2}{\mu-2L\delta}
\right],\nonumber\\
\Delta c
&=(\sqrt{a^+}+\sqrt h)^2-(\sqrt a+\sqrt h)^2.
\label{eq:certificate-perturbations}
\end{align}
For $h>0$,
\begin{align}
\Delta d
&\le(1+\sqrt{h/a^+})\Delta\beta
+\sqrt h\left(\frac1{\sqrt a}-\frac1{\sqrt{a^+}}\right)\beta\nonumber\\
&\quad+\frac{\sqrt{a^+}-\sqrt a}{\sqrt h}g,
\label{eq:certificate-delta-d}
\end{align}
and $\Delta d=\Delta\beta$ when $h=0$.
For $U(c,d)=d/(1-c)+c^n(R_0-d/(1-c))_+$, $|\partial_dU|\le(1-\bar c)^{-1}$ and $|\partial_cU|\le nR_0+\bar d(1-\bar c)^{-2}$.
The mean-value bound proves Eq.~\eqref{eq:robust-score-gap}; applying it to the upper- and exact-score minimizers gives $U_{\widehat a}\le U_{a^\star}+2\max_a\varepsilon_a$.

\subsection{Return transfer and arbitrary terminal indices}
\label{app:return}
Let $n=\lfloor N_a/(2D+1)\rfloor$.
The envelope argument keeps every later incomplete-block iterate below $\max\{R_{n(2D+1)},R_\star\}$, which is the truncated-transient form in Eq.~\eqref{eq:certificate-score}.
Since $J_\nu(w)-J_\nu(w^\star)=\psi_\nu^\top e$ for $\psi_\nu=\E_{s\sim\nu}\phi(s)$, Cauchy--Schwarz and the terminal envelope prove Corollary~\ref{cor:return}.

\subsection{Proof of Theorem~\ref{thm:tracking}}
\label{app:tracking}
Set $R_\star=d/(1-c)$ and $\Gamma_r=(Z_r-R_\star)_+$.
The explicitly assumed uniform-history drift and the positive-part replacement argument give
\begin{align}
\widetilde Z_{r+1}
&\le R_\star+c\Gamma_r,\nonumber\\
\Gamma_{r+1}
&\le(1+\chi)c\Gamma_r+\chi R_\star+(1+\chi^{-1})v^2.
\label{eq:tracking-old-target}
\end{align}
The second line follows by applying Young's inequality to the target shift and subtracting $R_\star$.
Iteration and $\Gamma_0\le Z_0$ prove Eq.~\eqref{eq:tracking-bound} whenever $\bar c=(1+\chi)c<1$.
Finally, $\widehat A_w-A^\star=\{\gamma\phi(s')-\phi(s)\}^\top(w-w_n^\star)$ and $\|\gamma\phi(s')-\phi(s)\|\le(1+\gamma)\Phi$; Cauchy--Schwarz proves Eq.~\eqref{eq:advantage-transfer}.

\subsection{Existence and complexity of the mixed controller}
\label{app:existence}
For fixed $(q,b)$, every coefficient and Eq.~\eqref{eq:certificate-score} is continuous on the declared positive-step domain.
The non-strict certification margins make that domain closed and bounded, so the extreme-value theorem applies whenever it is nonempty.
Finite enumeration across nonempty pairs gives a global minimizer.
A numerical implementation attains this oracle minimum when its scalar solver has certified global accuracy.
The $N_qN_b$ scalar searches cost $O(N_qN_bL_\eta)$ evaluations.
For rank-one $H_{i,k}$, feature inner products avoid dense matrices in update aggregation, which costs $O(qd)$ arithmetic and $O(d)$ memory.
Exact spectral moment calibration is charged separately from rank-one update aggregation.

\section{Correlation-Limited Minimax Risk}
For one round of the Gaussian location subclass, the agent average has variance $\sigma_c^2+\sigma_e^2/q$.
The $T$ round averages are independent and Gaussian with known variance, so their sample mean is minimax for squared loss over the unrestricted location family and has risk $T^{-1}(\sigma_c^2+\sigma_e^2/q)$.
Dividing the single-agent risk by this value gives Eq.~\eqref{eq:speedup}.
Since $1+(q-1)\rho\ge q\rho$ for $\rho>0$, the speedup is at most $\rho^{-1}$.

\section{Proof of the Stopped Adaptive Information Identity}
\label{app:information}
Condition on the pre-action history and realized action $A_k$.
An orthogonal rotation maps $X_k$ to $Y_k=q_k^{-1/2}\bm{1}^\top X_k$ and hypothesis-invariant standard-normal contrasts.
Because the predictable action kernels cancel, the stopped log likelihood is the sum of scalar Kalman innovation ratios, where $Y_k\sim\mathcal N(\sqrt{q_k}m^-_{j,k},V_{j,k})$.
The conditional Gaussian divergence and the bounded-stopping chain rule give Eq.~\eqref{eq:conditional-kl} and Eq.~\eqref{eq:adaptive-kl}; exchanging hypotheses gives the reverse identity.

For a decision $\widehat J$, set $p_v=\Pp_v\{\widehat J=j\}$.
Data processing gives
\begin{equation}
\KL(P_j^\pi|\F_\tau\,\|\,P_{1-j}^\pi|\F_\tau)\ge\operatorname{kl}(p_j,p_{1-j}).
\end{equation}
If both directional errors are at most $\delta$, monotonicity of binary relative entropy yields Eq.~\eqref{eq:info-constraint}.

\section{Occupation-Measure Opportunity Lower Bound}
\label{app:occupation}
Let $\Delta_j(z,a)\ge0$ be the one-step opportunity loss relative to the instance-$j$ certified oracle.
For a stopped predictable policy, define
\begin{equation}
\nu_j(G,a)=\E_j^\pi\sum_{k\le\tau}\mathbf{1}\{Z_k\in G,A_k=a\}.
\end{equation}
Every $\delta$-correct policy induces an occupation measure satisfying the controlled Kalman flow constraints,
\begin{align}
\int\mathcal{I}_{j,1-j}(z,a)\nu_j(dz,da)&\ge\operatorname{kl}(1-\delta,\delta),\nonumber\\
\int(h+q(a))\nu_j(dz,da)&\le B_{\rm msg},\nonumber\\
\int b(a)\nu_j(dz,da)+\ell\Pp_j(\tau>0)&\le B_{\rm env}.
\end{align}
Let $T_a$ be the belief transition and $\zeta_j$ the terminal belief law.
For every bounded measurable $f$, predictability and stopped telescoping give
\begin{align}
\int f(z)\zeta_j(dz)-f(z_1)
&=\int\left[\int f(z')T_a(z,dz')\right.\nonumber\\
&\left.\hspace{18mm}-f(z)\right]\nu_j(dz,da).
\label{eq:belief-flow}
\end{align}
Moreover,
\begin{equation}
\E_j^\pi\sum_{k\le\tau}\Delta_j(Z_k,A_k)
=\int\Delta_j(z,a)\nu_j(dz,da).
\end{equation}
Eq.~\eqref{eq:adaptive-kl}, Eq.~\eqref{eq:budgets}, and Eq.~\eqref{eq:belief-flow} place every admissible occupation measure in the controlled feasible set.
Taking the infimum of its exact loss over that larger set proves Proposition~\ref{prop:occupation}.

\section{Proof of the Main Learning-Value Guarantee}
\label{app:main}
Let $\mathcal H_d$ be the complete probe history.
The declared conditional certificate gives $\E_j[\|e_{\rm LA}\|^2\mid\mathcal H_d]\le\overline R_j(\widetilde a_{\widehat J};B_d^D,D)$ for either selected action.
The right side is at most $\overline R_j^{\rm sel}+\mathbf1\{\widehat J\ne j\}W_j(d)$.
Taking expectation and using the directional error bound proves Eq.~\eqref{eq:main-risk}.

Subtracting the full-budget all-agent certificate and using Eq.~\eqref{eq:value-decomposition} gives
\begin{equation}
\E_1\|e_{\rm LA}\|^2-\overline R_1^{\rm all}(B,D)
\le C_1^{\rm probe}+C_1^D-G_1+\delta W_1<0.
\end{equation}
The same decomposition on instance zero is bounded above by $\epsilon_{\rm safe}(d)\le\bar\epsilon$.
The fallback statement follows directly from the strict admission rule.

\section{Proof of the Threshold Sandwich}
\label{app:threshold}
For the independently initialized Gaussian running mean, let $c_a=\max\{(h+q)/\beta_m,b/\beta_e\}$ and $r=\lambda^b$.
The exact horizon is $N_a=s/c_a+O(1)$, with bounded terminal charges and rounding.
Eq.~\eqref{eq:mean-risk} and Eq.~\eqref{eq:geometric-correlation-sum} give $K_j(a)=c_a\{q^{-1}+\theta_j(1+r)/(1-r)\}$ and a uniform $O(s^{-2})$ remainder.
Removing probe charge $v$ changes $N_a$ by $O(v+1)$ and hence risk by $O((v+1)/s^2)$ for $s\ge s_0(v+1)$; compact persistence and the finite catalogue make these constants uniform.
Thus the phase-theorem risk conditions have a concrete Gaussian realization, separately from the constant-step affine certificate.
Write $d_\delta=\operatorname{kl}(1-\delta,\delta)$.
Write $r_0=1-\gamma<1$ and $Q=\max{\cal Q}$.
Under either true instance, the predictable spacing gives $\|C_{S_k}\|_{L^2}\le r_0\|C_{S_{k-1}}\|_{L^2}+\sqrt{\theta_+}$, hence a uniform bound $M_C=\sqrt{\theta_+}/(1-r_0)$.
The normalized observation $Y_k/\sqrt{q_k}$ has $L^2$ norm at most $M_C+1$.
For either candidate filter, write its posterior mean as $(1-t_{j,k})m^-_{j,k}+t_{j,k}Y_k/\sqrt{q_k}$, with predictable $0\le t_{j,k}\le1$.
Prediction multiplies this mean by a factor of magnitude at most $r_0$; the triangle inequality therefore bounds both predictive means uniformly in $L^2$ by $M_F=r_0(M_C+1)/(1-r_0)$, starting from zero.
Both predictive variances are at most $\theta_+$, so $1\le V_{j,k}\le V_+=1+Q\theta_+$.
Eq.~\eqref{eq:conditional-kl} now gives the expected, rather than pointwise, bound $\E_j\mathcal I_{j,1-j}(Z_k,A_k)\le i_+$, where $i_+=(\log V_++V_++4QM_F^2)/2$ is finite.
Extend a stopped rule predictably after stopping; nonnegativity and its deterministic budget bound $n_{\max}(s)\le c_Bs$ imply $\E_j\sum_{k\le\tau}\mathcal I_{j,1-j}\le i_+n_{\max}(s)$.
Together with Eq.~\eqref{eq:info-constraint}, this expected-information bound proves $B_N\ge c_Nd_\delta$ over Gaussian observation histories.
The fixed probe has information at least $i_->0$, so $n_\delta=\lceil i_-^{-1}\log(1/\delta)\rceil$ observations suffice and cost at most $c_pn_\delta+D_{\max}$.
Continuity of $\log(1/\delta)/d_\delta$ gives $n_\delta\le c_0+c_1d_\delta\le c_0+c_2B_N$.
Write $u_\delta\le c_pn_\delta+D_{\max}$ for the full probe charge.
The assumed expansion and budget-shift inequality give a high-instance improvement over the full-budget fallback certificate of at least $g_{\min}/s-\{2M+L(u_\delta+1)\}/s^2-\delta w_+/s$.
It is strictly positive for $s>\{4M+2L(u_\delta+1)\}/g_{\min}$, also requiring $s\ge s_0(u_\delta+1)$.
On instance zero, the assumed coefficient ordering gives $\epsilon_{\rm safe}\le\{2M+L(u_\delta+1)\}/s^2+\delta_0w_+/s$.
Thus $s\ge\sqrt{2\{2M+L(u_\delta+1)\}/\bar\epsilon}+2\delta_0w_+/\bar\epsilon$ suffices.
Using $2\sqrt{xy}\le x+y$, all three requirements are bounded by a class-constant multiple of $n_\delta+B_{\rm oracle}+B_\epsilon+1$.
Taking the maximum proves the upper bound in Eq.~\eqref{eq:threshold-sandwich}; its lower bound follows because every admitted $\delta$-correct probe obeys both information constraints.

